\documentclass[11pt]{article}
\usepackage[margin=0.85in]{geometry}
\usepackage[T1]{fontenc}
\usepackage{lmodern}
\usepackage{amsmath,amssymb,booktabs,array,tabularx}
\usepackage{xcolor,listings}
\usepackage[hidelinks]{hyperref}
\hypersetup{pdftitle={Assignment 2, Question 3(b): RTT--Loss AIMD},
  pdfauthor={CS422 Fall 2026 project team}}
\lstset{basicstyle=\ttfamily\small,columns=fullflexible,keepspaces=true,
  breaklines=true,frame=single,framesep=5pt,showstringspaces=false}
\setlength{\parindent}{0pt}
\setlength{\parskip}{0.5em}
\newcommand{\code}[1]{\texttt{\detokenize{#1}}}
\newcommand{\repo}{https://github.com/LelsersLasers/CS422-Fall2026}
\title{Assignment 2, Question 3(b)\\Design of RTT--Loss AIMD (RLA)}
\author{Ayusha Patra \and Kabileswaran Krishna Kumar \and
  Vihaan Kodeboyina \and Millan Kumar}
\date{CS 422: Computer Networks, Fall 2026}

\begin{document}
\maketitle

\section{Design objective and measurement evidence}

We propose \emph{RTT--Loss AIMD} (RLA), a window controller that combines
loss-triggered multiplicative decrease with a gentler response to RTT inflation.
Acknowledged throughput gates additive growth so that enlarging an unused window
is not the default response to a low sending rate. Standard slow start provides
initial capacity discovery. The objective is useful throughput with restrained
queueing, rather than maximum window growth regardless of delay.

The design is motivated by the current Assignment 2 comparison in
\code{output/part3/manifest.json} and \code{comparison_summary.csv}.
That manifest records Linux 6.18.50, 60-second transfers, one-second samples,
and the fixed order CUBIC, Reno, then BBR. There is \emph{one repetition}, with
10 CUBIC, 8 Reno, and 9 BBR successes across ten selected endpoints.
The examples below use only runs named by that summary; unrelated CSV files
left in the output directory are excluded.

\begin{table}[htbp]
\centering
\small
\begin{tabular}{llrr}
\toprule
Endpoint & Algorithm & Mean goodput (Mbit/s) & Median SRTT (ms) \\
\midrule
LeaseWeb Washington:5210 & CUBIC & 455.47 & 30.73 \\
 & Reno & 549.60 & 34.37 \\
 & BBR & 566.90 & 31.45 \\
moji Paris:5219 & CUBIC & 214.10 & 98.49 \\
 & Reno & 1.45 & 100.81 \\
 & BBR & 97.11 & 98.77 \\
\bottomrule
\end{tabular}
\caption{Motivating observations, not causal algorithm rankings. Mean goodput
is the arithmetic mean of interval rates; median SRTT is computed from each
run's TCP samples.}
\label{tab:evidence}
\end{table}

Washington refers to \code{speedtest.wdc2.us.leaseweb.net:5210}; its
CUBIC/Reno/BBR run identifiers are \code{d634f0b7e5ee},
\code{164436543996}, and \code{47b6fa078b58}.
Paris refers to \code{iperf3.moji.fr:5219}, with runs
\code{f392bf0bd8f0}, \code{db317c24e38b}, and \code{80dc7b613fb9}.
For example, the Washington traces show a throughput/delay tradeoff between
Reno and BBR. The Paris BBR trace includes a consecutive-sample increase of
1,741 in the cumulative retransmission counter. This motivates one reaction
per native recovery episode rather than one halving per retransmitted segment.
The very low Paris Reno rate, despite a similar median RTT, also cautions
against treating RTT alone as an explanation of throughput.

These sequential public-Internet runs do not establish equal path capacity,
background traffic, or fairness. They cannot identify an optimal threshold or
validate RLA. The missing repetitions required by Question 3(a) remain a
limitation of this evidence. We use the observations to motivate a testable
design, without claiming that RLA outperforms the measured algorithms.

\section{State, signals, and parameters}

\begin{center}
\small
\begin{tabularx}{\linewidth}{p{0.16\linewidth}X}
\toprule
Symbol & Meaning and units \\
\midrule
$W$, $M$ & Current congestion window in segments and sender MSS in bytes. \\
$R$ & Positive smoothed RTT in microseconds; Linux stores SRTT scaled by eight,
so the module divides \code{srtt_us} by eight. \\
$R_{\min}$ & Minimum positive SRTT observed by the controller over this connection. \\
$A$, $t$ & Cumulative acknowledged bytes and monotonic ACK-processing time in microseconds. \\
$W_e,M_e,R_e$ & Window, MSS, and SRTT captured at the start of the epoch. \\
$A_e,t_e,T_e$ & Epoch byte/time baselines and duration
$T_e=\max(100{,}000,R_e)$ microseconds. \\
$S$ & Slow-start threshold, in segments, maintained with Linux recovery. \\
\bottomrule
\end{tabularx}
\end{center}

After connection establishment, Linux supplies the initial window (normally
ten segments), threshold, MSS, and window clamp. RLA initializes its private
state to zero and learns the RTT baseline from positive SRTT observations.
The baseline is retained through recovery and idle restarts. An epoch is
initialized on the first valid, non-recovery ACK callback. No RTT sample,
zero MSS, or an invalid counter/time baseline authorizes window growth.

At the end of an epoch, let $\Delta t=t-t_e>0$ and $\Delta A=A-A_e$.
The acknowledged rate $G$ and utilization proxy $u$ are
\[
 G=\frac{8\Delta A\,10^6}{\Delta t}\quad[\mathrm{bit/s}],\qquad
 C_e=\frac{8W_eM_e\,10^6}{R_e}\quad[\mathrm{bit/s}],\qquad
 u=\frac{G}{C_e}=\frac{\Delta A R_e}{\Delta t W_eM_e}.
\]
$C_e$ is a rough window-supported rate, \emph{not} an estimate of bottleneck
capacity. The factor eight converts bytes to bits, and $10^6$ converts
microseconds to seconds. The implementation compares $5\Delta A R_e$ with
$4\Delta t W_eM_e$ using checked 64-bit products; an overflow disables
additive growth for that epoch rather than allowing arithmetic wraparound.

Relative RTT inflation is $q=(R-R_{\min})/R_{\min}$. RLA uses a low threshold
of 10\%, a high threshold of 25\%, a utilization threshold of 80\%, a loss
backoff factor of $0.5$, and a delay backoff factor of $0.9$.
The normal minimum is two segments, subject to the kernel's window clamp.
These are explicit starting heuristics, not fitted or validated constants.

For loss feedback, the implementation uses the kernel's
\code{ssthresh} callback rather than estimating independent
packet-loss probability from a coarse user-space sample. Fast retransmission,
RACK/SACK recovery, and timeouts remain Linux TCP responsibilities.
The same backoff is used for native ECN congestion responses.
The Python measurements separately log cumulative \code{total_retrans} as
a loss proxy and \code{lost} as the currently outstanding lost segments;
the latter is not cumulative. Neither counter is itself a packet-loss probability.

\section{Window-update pseudocode}

Loss handling and normal ACK processing are separate events. A loss response
has priority because Linux does not let normal RLA epoch decisions replace
its recovery processing. Reducing the window does not retract outstanding
bytes, and all sending remains subject to receiver flow control.

\clearpage
\begin{lstlisting}[basicstyle=\ttfamily\fontsize{10}{11.5}\selectfont]
INITIALIZE:
    keep Linux initial cwnd W, threshold S, and window clamp
    Rmin = unknown; epoch = invalid

ON NATIVE CONGESTION/LOSS THRESHOLD REQUEST:
    invalidate epoch
    return S = max(2, floor(W / 2))
    # Linux applies recovery/PRR or its RTO restart.

ON ENTERING OR LEAVING RECOVERY, OR NATIVE IDLE RESTART:
    invalidate epoch; retain Rmin

ON AN ACK ELIGIBLE FOR CONGESTION AVOIDANCE:
    if native recovery is active:
        invalidate epoch
        allow Linux/Reno-style timeout restart if in RTO loss state
        return                  # no custom delay or additive response

    read A, t, R, MSS, W, S, ACKed_segments, and cwnd_limited
    Rmin = min(Rmin, R) over positive R observations
    if R <= 0 or MSS <= 0 or W <= 0 or clamp <= 0:
        invalidate epoch; return
    if A or t went backwards:
        snapshot a new epoch; return
    if epoch is invalid:
        snapshot (Ae, te, We, Me, Re)
        Te = max(100000 microseconds, Re)

    due = (t - te >= Te)
    q = (R - Rmin) / Rmin
    if due and q > 0.25:
        W = min(clamp, max(2, floor(0.90 * W)))
        S = W                   # exit slow start as well
    else if W < S and cwnd_limited:
        W = min(clamp, S, W + ACKed_segments)  # standard slow start
    else if due and cwnd_limited and q < 0.10:
        if checked utilization calculation succeeds and u >= 0.80:
            W = min(clamp, W + 1)
    else:
        hold W

    if due:
        snapshot a new epoch using the resulting W
        # Exactly one epoch decision; no missed-epoch catch-up.
\end{lstlisting}

Only one branch changes the window on an eligible ACK. In the 10--25\%
delay band, congestion avoidance holds the window; startup can still increase
it until the high threshold or native threshold ends startup.
Threshold comparisons are strict for delay: exactly 10\% does not authorize
additive probing, and exactly 25\% does not trigger a delay decrease.
The 80\% utilization comparison is inclusive.

Normal loss backoff has a two-segment threshold, but native RTO handling can
temporarily set the actual window to one segment. RLA preserves this behavior.
It also retains Linux's undo behavior for spuriously detected congestion.
Recovery entry/exit and undo invalidate epoch baselines so that a native
decrease is not followed by a second decrease from the same mixed epoch.
The congestion controller does not halve once per duplicate ACK or per
retransmission counter increment.

\section{Why these rules}

\textbf{Slow start and shorter epochs.}
The earlier proposal started at ten segments and added only one segment per
second. With an MSS of 1,448 bytes and RTT of 100 ms, ten segments support
only about 1.16 Mbit/s; twenty support about 2.32 Mbit/s. That rule would
spend most of a ten-second experiment below the rates seen in our traces.
Standard slow start grows quickly while ACKs show the window is used.
Subsequent epochs are at least 100 ms and at least one starting SRTT,
giving the changed window time to affect feedback. This is faster than
one-second probing on ordinary measured RTTs, without introducing an
independently exponentially growing throughput target.

\textbf{Loss first, delay before more loss.}
The measured retransmission bursts motivate native episode-based recovery.
A 50\% threshold reduction is a conventional conservative response to
congestion feedback. A separate 10\% window reduction at elevated RTT can
drain a queue before loss occurs. The delay signal is relative inflation,
not a direct queue occupancy measurement. The hold band separates increase
and decrease thresholds to reduce switching from small RTT fluctuations.

\textbf{Throughput-gated additive probing.}
A large window does not guarantee goodput: an application, receiver, pacing,
or shared bottleneck can impose a lower rate. Requiring 80\% utilization
avoids automatically growing an apparently unused allowance. The kernel's
window-limited check supplies a second guard. These tests cannot distinguish
all causes of low throughput, and can conservatively stall an underserved
flow. A gate failure holds the window rather than increasing it.

\textbf{Relation to the bandwidth--delay product.}
At capacity $B$ bits/s and propagation RTT $R_p$ seconds, the propagation
BDP is $BR_p/8$ bytes. A roughly sufficient unqueued window is
$W\simeq BR_p/(8M)$ segments. Larger windows can build queues without
proportional goodput growth. RLA uses RTT inflation to discourage this
operating regime, but does not know $B$ or the true $R_p$.
For illustration, substituting the Washington BBR mean goodput
566.90 Mbit/s and its observed minimum SRTT of 28.041 ms gives about
1.99 MB. This is a \emph{measurement-based scale}, not a verified bottleneck
BDP: goodput may be below capacity and even minimum SRTT can contain queueing.

\section{Expected fairness and stability}

The fairness goal is similar long-run goodput for saturated RLA flows sharing
one bottleneck, initially assuming equal MSS and comparable propagation RTT.
If flows have synchronized epochs and common congestion feedback, equal
additive increments preserve their window difference, while a shared
multiplicative decrease scales that difference by 0.5 or 0.9. Repeated shared
decreases therefore tend to reduce window inequality. With equal RTT and MSS,
more equal windows can imply more equal rates.

Those conditions need not hold. Epochs are asynchronous, drops can affect
only some flows, and throughput gating can prevent the smaller flow from
growing. Epochs longer than 100 ms depend on RTT, so longer-RTT flows get
fewer update opportunities. Equal windows also give different rates at
different RTTs. Relative delay thresholds allow a larger absolute delay
increase on longer-baseline paths. Late arrivals may observe a persistently
queued RTT baseline. RLA can yield to Reno/CUBIC flows that keep growing
until loss, and no fairness guarantee is made against BBR or other controllers.
The available independent public-path measurements do not test shared-bottleneck
fairness.

Under stationary capacity, saturated use, and responsive feedback, additive
probing increases pressure at low delay, while high delay and native losses
reduce it. The hold band admits a possible operating region with
$0.10\le q\le0.25$ and no congestion events. This supports an expectation of
bounded operation or a sawtooth around an operating region, not a proof of
convergence to one fixed rate. Delayed feedback, simultaneous startup,
continuing losses, or a bandwidth change can sustain oscillations.
Conversely, low measured utilization can produce a stable but inefficient
hold. Stability, high utilization, and fairness must be evaluated separately.

\section{Assumptions, implementation, and limitations}

RLA assumes meaningful positive SRTT measurements, monotonic acknowledged-byte
counters, and a sender able to enforce a window through TCP. SRTT smooths
short events; ACK compression and cumulative ACK bursts distort epoch rates.
Changes in MSS or the window within an epoch make its utilization estimate
approximate, especially during startup. A lifetime RTT minimum can become
stale after a route change or can start too high behind a persistent queue.
Wireless losses, reordering, and spurious congestion indications can cause
unnecessary backoff despite Linux's recovery/undo mechanisms.

The 100-ms floor is deliberately conservative on short-RTT paths. Additive
recovery can still be slow on high-BDP paths after a large loss reduction.
Fast startup can overshoot before an epoch completes, and repeated delay
decreases can sacrifice throughput against more aggressive competing flows.
No finite thresholds establish universal fairness or convergence.

The \href{\repo}{project repository} contains the
\href{\repo/blob/millan/a2_3b_and_a3/Assignment3/kernel/tcp_rla.c}{RLA kernel module}.
This single source file implements window arithmetic, registers \code{rla},
supplies live TCP state, and delegates recovery to Linux.

Build and load commands are in \code{Assignment3/README.md}.
Integration with the Assignment 2 experiment pipeline, repeated four-algorithm
measurements, plots, and aggregate analysis are deferred to later work.
Repository links refer to the implementation branch and become available
when that branch is pushed.

Compared with the earlier \code{plans/part3b-design.md} proposal, the substantive
changes are standard slow start, $\max(100\,\mathrm{ms},R_e)$ epochs instead
of one-second epochs, epoch-start utilization snapshots, integer/clamp rules,
and native event-based recovery instead of periodic retransmission-based
halving. These changes make a concrete kernel implementation possible and
address startup underutilization. They have not established comparative
performance, fairness, or convergence.

\textbf{AI acknowledgment.} OpenAI Codex assisted with reading the assignments,
reviewing existing measurements, developing the revised algorithm and Linux
implementation and preparing this report section. The reported
measurements are identified by raw run files; hypotheses and limitations are
distinguished from demonstrated results.

\begin{thebibliography}{9}
\bibitem{assignment} CS 422, Fall 2026. \emph{Assignment 2}, Question 3(b),
and \emph{Assignment 3}, Option 1. Assignment PDFs supplied in this repository.
\bibitem{linux} Linux kernel sources, \code{net/ipv4/tcp_cong.c},
\code{net/ipv4/tcp_input.c}, and \code{include/net/tcp.h}.
\href{https://github.com/torvalds/linux/tree/v6.18/net/ipv4}{Linux v6.18 TCP source}.
The actual module was built against the running Linux 6.18.50 development tree.
\bibitem{aimd} M. Allman, V. Paxson, and E. Blanton.
\emph{TCP Congestion Control}, RFC 5681, 2009.
\url{https://www.rfc-editor.org/rfc/rfc5681}.
\bibitem{kbuild} Linux kernel documentation.
\emph{Building External Modules}.
\url{https://docs.kernel.org/kbuild/modules.html}.
\end{thebibliography}
\end{document}
